\section{Appendix C. The bridge between the two monetary references}\label{appendix-c.-the-bridge-between-the-two-monetary-references}

Write \(M_i=\int\widehat g_i\,d\nu\) for the mass of the opportunity object on the domain the reference integral uses, and \(o\) for the non-work bundle. The two references satisfy

\[
\beta_c\log(W^{1}_i/\lambda_c)=\log J_i-\log H_i,
\qquad
\beta_c\log(W^{4}_i/\lambda_c)=\log J_i-\log M_i-L_i(o),
\]

so that \(\log(W^{4}_i/W^{1}_i)=\big(\log H_i-\log M_i-L_i(o)\big)/\beta_c\) and \(\Delta_i=L_i(o)+\log M_i-\log H_i\). Omitting \(\log M_i\) is valid only when \(M_i=1\), and it is not: the audit reports median mass 0.077283 for single adults and 147.963 for couples. The premise that failed is therefore the unit mass of the kernel on that domain, not the maximality of the non-work bundle, which holds for every household in both samples, and not the orientation of the comparison.

After renormalizing the kernel on exactly that domain, the median gap is 0.0322 nats for single adults and 0.9906 for couples, the gap is non-negative for every household, and the ratio lies in \([0.9169, 0.9990]\) for single adults and \([0.5211, 0.7171]\) for couples. The identity above is satisfied to machine precision at the corrected normalization.

A separate point concerns the low-temperature limit. Attained and reference values must be indexed by the same shock scale for the limit to be meaningful. The statement that survives is that the two consistently indexed measures converge as the scale goes to zero; convergence to the value of staying at home evaluated at a different fixed scale is not the same statement and is not made.

